\chapter{Analytic approximation and divisorial bounds}
\label{chap:analytic-approximation}

Analytic approximation is a classical local and global construction on compact complex manifolds. Its availability should be separated from the convergence of Monge--Amp\`ere masses or of a metric on singularities. We recall the standard regularization theorem in the precise form needed for a nonclosed background, explain its normalization, and derive a uniform divisorial bound from Gauduchon metrics. The next chapter will apply this bound to restrictions of multiplier ideals.

Throughout, $X$ is connected and compact. We use the local Ohsawa--Takegoshi theorem, coherence of multiplier ideals, and Demailly's classical global regularization theorem as standard reference inputs. They are stated below in the forms used. The global regularization input is \cite[Theorem 14.12 and its construction]{DemAM11}; the numbered locations refer to the author's July 2011 text. It permits an arbitrary continuous lower bound, without requiring that bound to be closed.

\section{Analytic types and restriction}

We follow the convention of \cite{SGPT}: an analytic singularity has, locally, a presentation
\[
 u=c\log\bigl(|f_1|^2+\cdots+|f_N|^2\bigr)+O(1),
 \qquad c\in\mathbb Q_{>0},
\]
with holomorphic generators not all identically zero. The remainder is bounded; smoothness of this remainder is an additional property. A coherent ideal $\mathfrak a$ and a common coefficient $c$ specify the \emph{analytic type} $(c,\mathfrak a)$. Changing local generators changes the logarithm by a bounded amount on smaller relatively compact open sets.

\begin{proposition}\label{prop:analytic-type-calculus}
Analytic singularities are preserved by addition, finite maxima, and multiplication by a positive rational number. If $Y\subset X$ is a connected complex submanifold and the ideal of an analytic presentation does not vanish identically on $Y$, the restriction is genuine and has analytic type
\[
 (c,\mathfrak a\mathcal O_Y).
\]
Here $\mathfrak a\mathcal O_Y$ means the image ideal generated by the restrictions of local sections of $\mathfrak a$.
\end{proposition}
\begin{proof}
For a rational scaling, multiply the coefficient. For two presentations, choose a common positive rational coefficient $d$ with $c_1=k d,c_2=\ell d$, where $k,\ell$ are positive integers. Products of generators present the sum, since a product of sums of squares is a sum of squares of pairwise products. For the maximum use
\[
 \max\{A,B\}\leq A+B\leq2\max\{A,B\},\qquad A,B\geq0,
\]
applied to the appropriate sums of squares of generators of $\mathfrak a_1^k$ and $\mathfrak a_2^\ell$. Taking logarithms gives the required bounded error. These are local assertions; finite covers give the corresponding ordinary singularity comparisons on a compact manifold.

For restriction, restrict the displayed inequalities defining the bounded remainder. At least one generator remains nonzero on the connected manifold $Y$, so its common zero set is proper. The restricted function is quasi-plurisubharmonic, is not identically $-\infty$, and has the stated presentation. A smooth local correction gives the same conclusion directly from the usual restriction theorem for psh functions.
\end{proof}

Zero generic Lelong number does not make the restriction of an arbitrary quasi-plurisubharmonic function genuine. For example, on a small polydisc with $|z_1|<\mathrm{e}^{-1}$,
\[
 u(z)=-\sqrt{-\log|z_1|^2}
\]
is plurisubharmonic: the function $t\mapsto-\sqrt{-t}$ is increasing and convex for $t<0$. Along $\{z_1=0\}$, its Lelong number is zero but its restriction is identically $-\infty$. The analytic hypothesis in \cref{prop:analytic-type-calculus} will exclude this behavior for the approximants below.

\section{The local Bergman estimates}

Let $\psi$ be psh on a bounded pseudoconvex domain $\Omega\subset\C^n$. For a positive integer $m$, let
\[
 \mathcal H_m=\left\{f\in\mathcal O(\Omega):
          \|f\|_m^2:=\int_\Omega|f|^2\mathrm{e}^{-m\psi}\,\mathrm{d}\lambda<\infty\right\},
 \qquad
 \beta_m=\frac1m\log\sum_j|g_{m,j}|^2,
\]
where $(g_{m,j})_j$ is an orthonormal basis. The sum is equivalently the squared norm of point evaluation:
\[
 \sum_j|g_{m,j}(x)|^2
   =\sup_{\|f\|_m\leq1}|f(x)|^2.
\]
It is independent of the basis. The standard local finite-generation part of Bergman regularization says that $\beta_m$ has analytic type
$(1/m,\mathcal I(m\psi))$ on smaller open sets. This follows from the weighted holomorphic generation of the multiplier ideal and local $L^2$ division; see \cite[Proposition 5.7 and Section 14.A]{DemAM11}. We use that local analytic statement as a standard input, together with the point-extension estimate used in the proof below.

\begin{proposition}\label{prop:bergman-local-estimates}
There are constants $C_1,C_2$ independent of $m$ such that, for $B(x,r)\Subset\Omega$,
\begin{equation}\label{eq:bergman-local-estimates}
 \psi(x)-\frac{C_1}{m}
 \leq\beta_m(x)
 \leq\sup_{B(x,r)}\psi+\frac1m\log\frac{C_2}{r^{2n}}.
\end{equation}
Consequently $\beta_m\to\psi$ pointwise and in $L^1_{\mathrm{loc}}$, and
\begin{equation}\label{eq:bergman-lelong-bounds}
 \max\{\nu(\psi,x)-n/m,0\}
 \leq\nu(\beta_m,x)\leq\nu(\psi,x).
\end{equation}
\end{proposition}
\begin{proof}
For $\|f\|_m\leq1$, the submean inequality gives
\[
 |f(x)|^2
 \leq C_2r^{-2n}\int_{B(x,r)}|f|^2\,\mathrm{d}\lambda
 \leq C_2r^{-2n}\mathrm{e}^{m\sup_{B(x,r)}\psi}.
\]
Taking the supremum over the unit ball and then the logarithm proves the upper bound.

The local Ohsawa--Takegoshi point-extension theorem gives a holomorphic $f$ with $f(x)=1$ and
\[
 \int_\Omega|f|^2\mathrm{e}^{-m\psi}\,\mathrm{d}\lambda
 \leq C \mathrm{e}^{-m\psi(x)}
\]
when $\psi(x)>-\infty$. The constant is independent of the psh weight and of $m$. Normalize $f$ in $\mathcal H_m$. Point evaluation is then at least $C^{-1}\mathrm{e}^{m\psi(x)}$, proving the lower bound. At a pole the lower bound has no content and remains valid.

Choosing $r=1/m$ proves pointwise convergence by upper semicontinuity of $\psi$. On a fixed relatively compact set, the same estimates bound $\beta_m$ above by a constant and below by $\psi-C_1/m$. Since $\psi$ is locally integrable, dominated convergence also gives convergence in $L^1_{\mathrm{loc}}$.

For the Lelong estimate, take suprema over $B(x,r)$ in the two inequalities and bound the upper one using $B(x,2r)$. Divide by $\log r^2<0$ and let $r\downarrow0$. The term $-(2n/m)\log r$ contributes $-n/m$ after the division. The remaining constant terms disappear. This proves the two inequalities in \eqref{eq:bergman-lelong-bounds}; nonnegativity supplies the maximum with zero.
\end{proof}

In \cite{DemAM11}, the local weight is $\mathrm{e}^{-2m\phi}$ and Lelong numbers are normalized using $\log|z|$. Putting $\psi=2\phi$ gives exactly the formulas above. In particular the error is $n/m$ with the present weight $\mathrm{e}^{-m\psi}$ and normalization by $\log|z|^2$.

The same local construction and estimates hold for every real parameter $m>0$. Its local logarithmic presentation has coefficient $1/m$; this coefficient is rational exactly when $m$ is rational. We use the term analytic type only with the rational-coefficient convention above. The global sequence below is indexed by integers.

\Needspace{23\baselineskip}
\section{The global curvature bound}

We record precisely the classical global input, including the ideal type supplied by its construction.

\begin{theorem}[Demailly regularization]\label{thm:demailly-hermitian-input}
Let $u$ be quasi-plurisubharmonic on a connected compact complex manifold and let $\gamma$ be a continuous real $(1,1)$-form with $\ddc u\geq\gamma$. Fix a Hermitian form $\omega$. There are quasi-plurisubharmonic $u_m$, for sufficiently large positive integers $m$, and numbers $\epsilon_m\downarrow0$ such that
\begin{enumerate}
\item $u_m$ has analytic type $(1/m,\mathcal I(mu))$;
\item $u_m\geq u$ and $u_m\to u$ pointwise and in $L^1(X)$;
\item $\ddc u_m\geq\gamma-\epsilon_m\omega$;
\item $\max\{\nu(u,x)-n/m,0\}\leq\nu(u_m,x)\leq\nu(u,x)$ for all $x$.
\end{enumerate}
No monotonicity of the values $u_m$ is asserted here.
\end{theorem}

This is the form of \cite[Theorem 14.12]{DemAM11} used in this book. Its local generators are the weighted Bergman generators of $\mathcal I(mu)$; smooth local corrections do not change that ideal. The overlapping sums have comparable norms, so the glued logarithm retains this ordinary analytic type. The global construction freezes the continuous lower bound on sufficiently small coordinate balls. Its curvature error combines the modulus of continuity of that bound with an overlap term tending to zero as the Bergman index increases. It is this global estimate, rather than merely the local convergence in \cref{prop:bergman-local-estimates}, which is the standard input in the theorem.

\begin{corollary}\label{cor:nonclosed-analytic-approximation}
For any smooth real $\theta$ and $u\in\PSH(X,\theta)$, the approximants can be chosen so that
\[
 \theta+\ddc u_m\geq-\epsilon_m\omega.
\]
Their analytic type is $(1/m,\mathcal I(mu))$.
\end{corollary}
\begin{proof}
Apply the theorem to the closed, almost positive current $\ddc u$ with lower bound $\gamma=-\theta$. The theorem requires no closedness of $\gamma$. Adding $\theta$ gives the asserted inequality. In the convention of \cite{DemAM11}, this application is to the potential $u/2$, since $\frac{\mathrm{i}}{\pi}\partial\bar\partial(u/2)=\ddc u$.
\end{proof}

Thus the nonclosed background does not prevent analytic regularization with a vanishing curvature error. The issue in an $\mathcal I$-goodness theorem is what further topology or mass identity that approximation satisfies. Taking ordinary envelopes along dyadic indices does give a canonical limit: \cref{thm:dyadic-analytic-envelope-limit} identifies it with $Q_\theta[u]$. In a Hermitian background, rational scaling absorbs the curvature errors and gives analytic approximants in the original background, by \cref{cor:fixed-hermitian-analytic-density}. These conclusions concern singularities and envelopes; they do not prove convergence of the unrestricted Monge--Amp\`ere masses.

\Needspace{16\baselineskip}
\section{Gauduchon metrics bound divisorial orders}

A linear mass estimate suffices for a basic uniform bound on singularities. It is available on every compact Hermitian manifold.

\begin{lemma}\label{lem:gauduchon-uniform-lelong}
Let $Z$ be connected compact complex and let $\beta$ be a smooth real $(1,1)$-form with $\PSH(Z,\beta)\ne\varnothing$. There is a finite constant $C_\beta$ such that
\[
 \nu(v,z)\leq C_\beta\qquad
 \text{for every }v\in\PSH(Z,\beta)\text{ and every }z\in Z.
\]
In particular all generic divisorial Lelong numbers in this fixed background are uniformly bounded. Additive normalization of the potentials is unnecessary.
\end{lemma}
\begin{proof}
Write $d=\dim Z$. For $d\geq2$, choose a Gauduchon metric $\eta$. For $d=1$, use $\eta^{d-1}=1$. The conserved linear trace in \cref{prop:gauduchon-trace} gives
\[
 \int_Z(\beta+\ddc v)\wedge\eta^{d-1}
 =\int_Z\beta\wedge\eta^{d-1}.
\]
The current on the left is positive, so this fixed nonnegative number bounds its mass on every coordinate ball.

On a slightly larger ball choose a smooth quadratic $r$ with $\ddc r\geq\beta$. Then
\[
 T=\ddc(v+r)=(\beta+\ddc v)+(\ddc r-\beta)
\]
is closed and positive. Equivalence of fixed smooth metrics bounds its local trace mass by a constant times the preceding global trace plus a fixed smooth contribution. The local Lelong monotonicity formula bounds the Lelong number at the center by the trace mass on a fixed surrounding ball, with a constant depending only on that ball. Use finitely many nested coordinate balls to obtain the same bound for every center in the smaller balls. Smooth corrections do not change Lelong numbers. This proves the assertion.
\end{proof}

\begin{corollary}\label{cor:restricted-analytic-divisorial-bound}
Let $Y\subset X$ be a connected compact complex submanifold and suppose
\[
 \nu(u,Y):=\inf_{y\in Y}\nu(u,y)=0.
\]
For the approximants of \cref{cor:nonclosed-analytic-approximation}, each $u_m|_Y$ is genuine. If $\pi:Z\to Y$ is a smooth compact modification and $E\subset Z$ is a prime divisor, then
\[
 \sup_m\nu\bigl((u_m|_Y)\circ\pi,E\bigr)<\infty.
\]
\end{corollary}
\begin{proof}
For each $m$, choose $y\in Y$ with $\nu(u,y)<1/m$. The local Skoda integrability criterion gives $\mathcal I(mu)_y=\mathcal O_{X,y}$. Thus the image ideal $\mathcal I(mu)\mathcal O_Y$ is not identically zero. Analytic restriction in \cref{prop:analytic-type-calculus} makes $u_m|_Y$ genuine. Its pullback to $Z$ is also genuine.

\Needspace{6\baselineskip}
The restricted and pulled-back positivity inequality gives
\[
 \pi^*(\theta|_Y)+\ddc\bigl((u_m|_Y)\circ\pi\bigr)
 \geq-\epsilon_m\pi^*(\omega|_Y).
\]
A single smooth Hermitian form on $Z$, multiplied by a large enough constant, dominates these lower bounds for all indices. The preceding lemma therefore applies on $Z$ and bounds the generic Lelong number along $E$.
\end{proof}

This use of Gauduchon metrics controls a degree-one current after restriction and modification. It makes no claim about the top-degree masses of the restricted approximants. In particular it remains available when the ambient bounded mass property fails.
